\section{Related Work}

\subsection{Attacks on Tool-Using and Web Agents}
Research on language-model security began with adversarial inputs to a single model but has moved decisively toward tool-mediated, multi-step settings. Indirect prompt injection showed that content fetched by an application can steer the model that reads it~\cite{greshake2023what,liu2023prompt}, and automated search over injection payloads generalizes the attack~\cite{liu2024automatic}. When the application is an agent, the injected instruction becomes executable: InjecAgent measures direct and indirect injections against tool-integrated agents~\cite{zhan2024injecagent}, AgentDojo provides a dynamic environment with adaptive attackers~\cite{debenedetti2024agentdojo}, and the Agent Security Bench formalizes attack scenarios across agent scenarios~\cite{zhang2024agent}. Web-specific work shows that hostile page content can hijack browsing agents to leak private data~\cite{liao2024environmental} or manipulate multimodal action generation~\cite{wang2025webinject}; protocol-level studies extend the surface to tool discovery and selection under MCP~\cite{hou2025model,zhang2025security,wang2026adaptools}; and red-teaming frameworks adaptively search web agents for injection chains~\cite{syros2026muzzle}. Security analyses attribute the elevated risk of web agents to multi-step processing of untrusted content~\cite{chiang2025agents}, and safety benchmarks for autonomous web agents document policy-violating completions~\cite{tur2025safearena,levy2024stwebagentbench}. Recent measurements emphasize a property directly relevant to detection: attacks can be made \emph{covert}, leaving no user-visible trace of the compromised operation~\cite{lee2026will}. Our work takes this multi-step, state-changing threat model as given and asks what representation makes the resulting risk detectable and attributable.

\subsection{Defenses: Content Hardening, Guards, Judges, and Runtime Monitors}
Defenses cluster by the object they inspect. Content-level methods restructure queries or align preferences so instructions hidden in data are not followed~\cite{chen2024struq,chen2024secalign}; attention-based detectors localize anomalous attention mass over injected spans~\cite{zhong2025attention}; provable and design-level approaches enforce data-flow and capability constraints that make injections inert~\cite{zhu2025melon,debenedetti2025defeating}. Guard models classify inputs and outputs against safety taxonomies~\cite{inan2023llama}, with known calibration limitations~\cite{liu2024calibration}. Judge-style assessment scores whole interactions or trajectories~\cite{zheng2023judging,yuan2024rjudge,zhuge2024agentasajudge}, and analysis of outcome-only trajectory judges shows they miss faults that do not alter the final result, arguing for process-level evidence~\cite{mohammadi2026trajectoryjudge}. Runtime monitors predict risky next actions or reason over world models~\cite{wang2025probguard,lin2026dreamguard}, adapt lifelong~\cite{luo2025agrail}, operate at step level~\cite{mou2026toolsafe}, or produce diagnostic explanations~\cite{liu2026agentdog}; guardrails themselves, however, can become denial-of-service targets and their verdicts are formulation-sensitive~\cite{zhou2026from}. Across these families, outputs are either a scalar over text or a verdict over a run; the operation-level evidence needed to reconfigure an agent is rarely produced, and when it is (as in diagnostic frameworks) it is not coupled to a learned detector that can rank \emph{which} operations are suspect. That coupling---joint detection, localization, and conversion into agent-consumable knowledge---is the gap this paper addresses.

\subsection{Graph Structure, Graph Anomaly Detection, and MIL}
Modeling reasoning as a graph has precedent in prompting frameworks that branch and merge thoughts~\cite{yao2023tree,besta2023graph}, and dependency-guided attribution over execution graphs recently showed that decisive failures must be located with structural context~\cite{deng2026defa}. On the modeling side, we build on attention-based and relational graph encoders~\cite{velikovi2017graph,brody2021attentive,schlichtkrull2017modeling,hu2020heterogeneous}. Graph anomaly detection offers detectors at node and graph levels~\cite{tang2022rethinking,liu2021anomaly,roy2024gadnr,duan2022arise}, including substructure-aware~\cite{duan2022arise} and subgraph-aligned formulations~\cite{hu2025samcl}, anomaly detection over edge-attributed multigraphs~\cite{sotiropoulos2023adamm}, and transferable graph-level scoring~\cite{wang2026glass}. These methods, however, target social, transaction, or sensor graphs whose anomalies are unusual nodes or subgraph motifs; none instantiates a graph over agent executions, where the anomaly of interest is a \emph{semantic} dependency pattern---an operation and its justification chain---rather than a structural outlier, and none provides operation-anchored localization. Multiple instance learning supplies the weak-supervision discipline we need: bag-level labels with instance-level ambiguity~\cite{carbonneau2016multiple}, attention-based aggregation~\cite{ilse2018attentionbased}, and graph-based bag encoders~\cite{tu2019multiple}. We instantiate bags as trajectories, instances as \aeu{}-anchored \csg{}s, and add paired supervision that classical MIL does not use.

\subsection{Experience, Reflection, and the Closest Work}
Agents already convert experience into reusable knowledge: verbal self-reflection improves retries~\cite{shinn2023reflexion}, experience-driven learners extract and transfer insights across tasks~\cite{zhao2023expel}, and recent systems industrialize reflective learning~\cite{allard2026experiential}; a unified review describes this trend as \emph{externalization} of memory, skills, and protocols around the model~\cite{zhou2026externalization}. This literature reuses experience for \emph{competence}; we reuse it for \emph{safety}, grounding each experience in a localized risk subgraph rather than a free-text lesson. The closest work to ours is DeFA~\cite{deng2026defa}, which attributes failures to decisive errors using dependency guidance over agent traces; DeFA is post-hoc---it explains failures after they are known---whereas our formulation detects anomalous trajectories \emph{before} outcomes are trusted, localizes risk at the granularity of complete tool operations under only graph-level supervision, and additionally converts the localized evidence into \sexp{}s that condition future behavior. Runtime governance proposals argue for contracts and policies at execution time~\cite{ng2026agent,joshi2026deontic,errico2025securing}; our detector and experience format are designed as the concrete substrate such contracts would audit.
